%% Section de définitions — exigence d'autonomie du document.
%% Toute notation et tout terme interne employés dans la note sont ici.
\section{Définitions}
\label{sec:definitions}

Un lecteur qui n'a jamais vu ce projet doit pouvoir lire la note sans rien
chercher ailleurs. Toutes les notations et tous les termes internes employés
dans la suite sont donc rassemblés ici.

\subsection{Paramètres du modèle}

\begin{longtable}{@{}p{0.10\textwidth}p{0.86\textwidth}@{}}
\toprule
\textbf{symbole} & \textbf{sens, et plage explorée} \\
\midrule
\endhead
$\lambda$ & \textbf{intensité de naissance} : nombre moyen d'entités créées
  par pas, paramètre d'une loi de Poisson. Fixe la \emph{taille} du système
  sans changer ses grandeurs intensives. Exploré sur $[1\,;100]$. \\
$K_0$ & \textbf{dotation de naissance} : capital reçu par une entrante, en
  joules. Protège les jeunes. Exploré sur $[2\,;200]$. \\
$\sigma$ & \textbf{volatilité individuelle} : écart-type du choc log-normal
  multiplicatif appliqué au capital à chaque pas. Levier dominant de la
  démographie. Exploré sur $[0\,;1]$. \\
$\delta$ & \textbf{taux de dépréciation} du capital par pas. Fixe l'échelle
  du capital. Exploré sur $[0{,}01\,;0{,}10]$. \\
$k$ & \textbf{taille du bassin de rencontre} : nombre d'entités tirées à
  chaque tour de marché, parmi lesquelles la plus riche prête à la plus
  pauvre. Exploré sur $\{2,\ldots,10\}$ dans M4B ; vaut 2 et devient un choix
  constitutif dans les lignées récentes. \\
$A$ & \textbf{productivité} de la technologie, facteur du terme de
  production. C'est le levier d'« amélioration d'efficacité » de tout le volet
  rebond. \\
$\gamma$ & \textbf{exposant de production}, $0 < \gamma < 1$ : la production
  $A K^\gamma$ est \emph{concave} en la taille. Vaut $1/2$ dans la lignée de
  référence. \\
$\rho$ & \textbf{intensité de la phase de marché} : nombre de tours de
  rencontre par entité et par pas. Vaut 1 par défaut. C'est le seul levier
  qui déplace les exposants de queue (§\ref{sec:controle}). \\
$p$ & \textbf{partage du surplus coopératif}, imposé dans $[0,1]$ pour
  \code{bargain}, mais potentiellement hors de cet intervalle pour une autre
  règle : part du surplus
  de la paire que le service d'intérêt transfère à la prêteuse. $p = 0$
  altruiste, $p = 1$ asservissement total (§\ref{sec:partage}). \\
\bottomrule
\end{longtable}

\subsection{Variables d'état et observables}

\begin{longtable}{@{}p{0.15\textwidth}p{0.81\textwidth}@{}}
\toprule
\textbf{notation} & \textbf{définition} \\
\midrule
\endhead
$K_i$ & \textbf{capital}, ou taille intrinsèque productive de l'entité $i$, en
  joules. Réel, productif, dépréciable. \\
$C_i$, $D_i$ & \textbf{créances} et \textbf{dettes} de $i$ : sommes des
  principaux prêtés et empruntés. Nominaux, perpétuels, non amortis. \\
$NW_i$ & \textbf{valeur nette}, $NW_i = K_i + C_i - D_i$. C'est le prédicat de
  faillite : $NW_i < 0$ met l'entité en défaut. \\
$q$, $r$ & \textbf{principal} et \textbf{taux par pas} d'un contrat. Le service
  versé par pas vaut $rq$. \\
$L$, $\Delta$ & pour un contrat : \textbf{perte de puissance extractrice} de la
  donneuse, et \textbf{surplus coopératif} de la paire --- le gain de production
  jointe permis par le transfert de stock. Le service se décompose en
  $rq = L + p\Delta$. \\
\code{prod} & \textbf{production} d'une entité sur un pas, $A K_i^\gamma$. Le
  terme « puissance extractrice » désigne la même quantité. \\
\code{int\_in}, \code{int\_out} & \textbf{intérêts reçus} et \textbf{versés} par
  pas. \code{int\_in} est, dans l'ontologie retenue, l'analogue du revenu
  monétaire usuel. \\
$E_t$ & \textbf{énergie totale} présente dans le système en fin de pas,
  $E_t = \sum_i K_i(t)$. Distincte de la production totale : la concavité de la
  production et les variations d'effectif changent leur relation. \\
$\varepsilon$ & \textbf{élasticité} de la production totale à la productivité.
  La dérivée locale est $\mathrm{d}\ln(\mathrm{prod})/\mathrm{d}\ln A$ ;
  M4.4 estime un contraste fini : moyenne par graine de
  $\ln(\overline{Y}_{\text{traité}}/\overline{Y}_{\text{contrôle}})/\ln(1{,}5)$,
  sur $3000<t\le4000$. \\
$E$ & \textbf{réponse normalisée} : réponse observée divisée par la réponse
  strictement proportionnelle $(m-1)p_{\mathrm{trait}}$, où $m$ est le
  multiplicateur appliqué à $A$ et $p_{\mathrm{trait}}$ la part de production
  traitée avant intervention. $E > 1$ est qualifié de rebond dans le modèle. \\
\bottomrule
\end{longtable}

\subsection{Termes internes}

\begin{description}[leftmargin=1.3em,style=nextline]
\item[Bassin] L'échantillon de $k$ entités tiré à un tour de marché. Le
  \emph{Gini du bassin} est le coefficient de Gini du capital mesuré sur les
  entités entrant en rencontre, avant appariement.
\item[Carnet] L'ensemble des contrats de prêt vivants à un instant donné. Le
  \emph{taux moyen du carnet} est la moyenne des $r$ des contrats vivants. Un
  contrat quitte le carnet dès que l'une de ses deux extrémités meurt.
\item[Rotation du crédit] Volume de nouveaux contrats conclus par pas, rapporté
  au capital total. Mesure la vitesse à laquelle le stock productif change de
  mains. Ce n'est pas un nombre de contrats : c'est un flux rapporté à un stock.
\item[Tension d'échelle] Rapport $T = K_{\mathrm{aut}}/K_{\mathrm{eq}}$ entre le
  point fixe individuel sans crédit et le capital équivalent du système. Il
  résume bien les effets d'échelle \textbf{à $\delta$ fixé}, et \emph{échoue}
  comme paramètre d'état quand $\delta$ varie\footnote{\incertitude{} Le mot
  « tension » a été employé dans deux sens dans les documents antérieurs :
  celui-ci, et par extension la chaîne causale complète de la contraction de
  population. Dans cette note il désigne \emph{uniquement} le rapport ci-dessus,
  avec sa réserve ; la chaîne causale est nommée « la chaîne »
  (§\ref{sec:chaine}). Cette distinction terminologique est retenue dans
  la présente rédaction.}.
\item[Capital équivalent $K_{\mathrm{eq}}$] Capital qui, s'il était détenu par
  chaque entité, donnerait la production totale observée. Il sépare ce qui
  relève de l'effectif de ce qui relève de l'échelle par entité.
\item[Avalanche] Composante faiblement connexe du graphe des pertes de créances,
  restreinte aux mortes d'une même résolution de cascade
  (§\ref{sec:avalanches}).
\item[Racine] Faillite entrée directement dans la file initiale d'un pas, sans
  avoir été déclenchée par la perte d'une créance.
\item[Rapport de branchement $b_1$] $1 - \text{racines}/\text{morts}$ : fraction
  moyenne de morts qui ne sont pas des racines. Borné par 1 \emph{par
  construction}.
\item[Descendance moyenne $b_2$] Nombre moyen de couples causals par mort, lu
  sur le graphe des cascades. Peut dépasser 1 quand une victime a plusieurs
  parents.
\item[Sur-détermination] Distinguer trois mesures. L'excès d'arêtes par
  morte est $b_2-b_1=\sum_{v\text{ non racine}}(d_v-1)/N$, où $d_v$
  est le nombre de parents retenus à la génération précédente et $N$ le
  nombre total de mortes. Ce n'est pas la proportion de victimes à plusieurs
  parents parmi les non-racines vues (\code{multi\_parent\_share}) :
  une victime à trois parents contribue deux arêtes supplémentaires.
  La suffisance d'une perte pour déclencher une faillite est encore une autre
  propriété, étudiée par le plancher de suffisance du lot J.
\item[Fardeau] Somme des principaux dus par une entité, rapportée à sa
  production. C'est la variable dans laquelle la convexité du risque était
  posée en hypothèse (§\ref{sec:partage}).
\item[Contraste apparié] Comparaison graine par graine entre un bras traité et
  le contrôle qui partage son amorçage et l'état du générateur à la bifurcation.
  Cela ne garantit pas des chocs communs par entité après divergence des
  effectifs. Les intervalles des contrastes moyens de M4.4 sont des
  intervalles de Student à onze degrés de liberté ; les tableaux de dispersion
  des exposants emploient aussi des écarts-types, explicitement distingués.
\item[Bras] Une condition expérimentale d'une campagne : un jeu de paramètres
  appliqué à l'ensemble des graines.
\item[Portée] Étendue d'une intervention : \emph{globale} (toutes les entités),
  \emph{nouvelles} (les entrantes seulement) ou \emph{fraction} (une part des
  vivantes).
\item[Amorçage (\emph{burn-in})] Pas initiaux écartés de toute mesure, le temps
  que le système atteigne son régime stationnaire.
\item[Élasticité] $\mathrm{d}\ln Y / \mathrm{d}\ln X$ : variation relative de
  $Y$ par variation relative de $X$. Sauf mention contraire, $X$ est la
  productivité $A$.
\item[Coefficient de Gini] Mesure d'inégalité entre 0 (tous égaux) et 1 (tout
  chez un seul). Calculé ici sur le capital, la valeur nette, la production ou
  le revenu d'intérêt --- quatre grandeurs qu'il ne faut jamais confondre,
  puisqu'elles donnent des valeurs allant de $0{,}05$ à $0{,}61$ sur le même
  état.
\item[PRCC] Coefficient de corrélation partielle de rang, mesuré sur un plan
  d'expérience par hypercube latin. Entre $-1$ et $1$ ; mesure l'association
  monotone d'un paramètre à une sortie, les autres tenus constants.
\item[Statistique de Vuong] Test de comparaison de deux lois non emboîtées par
  rapport de vraisemblance. $|z| > 1{,}96$ tranche en faveur de l'une ; en deçà
  le test ne conclut pas --- ce qui ne signifie pas que les deux lois s'ajustent
  aussi bien, mais que \emph{ce test} ne les sépare pas.
\item[Estimateur de Hill] Estimateur de l'exposant de queue d'une loi de
  puissance, calculé sur les points au-delà d'un seuil. La \emph{couverture} est
  le nombre de points retenus au-delà de ce seuil ; elle gouverne la précision.
\item[CCDF] Fonction de survie empirique, $P(S \ge s)$. Tracée en échelle
  logarithmique, une loi de puissance y est une droite et une coupure de taille
  finie un décrochement.
\end{description}
